\begin{abstract}
LLM-based automated program repair requires localizing relevant code within large repositories, yet the dominant retrieval-augmented generation (RAG) paradigm---embedding, indexing, and querying repository contents---incurs substantial cost that is rarely reported alongside resolve rates. We propose \emph{Graph-Map Progressive} (GM-P), a lightweight context-construction strategy that builds a language-level dependency graph from the repository and progressively expands context from the suspected fault locale along structural edges. By replacing embedding-based retrieval with deterministic graph traversal, GM-P eliminates per-query vector-search costs while preserving the LLM's ability to reason over cross-file dependencies. On SWE-bench Verified (500 instances across 12 repositories), GM-P resolves 46.8\% of issues---significantly more than every compared method after Holm--Bonferroni correction of paired McNemar tests ($\alpha=0.05$)---against an oracle upper bound of 53.8\%. Crucially, GM-P consumes approximately 256K tokens per resolved issue, roughly 7.5$\times$ fewer than the nearest RAG-based progressive method (1.93M tokens), because the graph is built once with deterministic tooling and amortized across all issues within a repository snapshot.\footnote{Both graph and RAG indices were constructed in our experimental harness; reported costs account for method-specific token consumption only. See Section~\ref{sec:threats} for discussion.} These results suggest that structural dependency graphs offer a cost-effective and statistically reliable alternative to embedding-based retrieval for repository-level issue localization and code repair.
\end{abstract}
